\documentclass[10pt,a4paper]{amsart}
\usepackage[margin=19mm]{geometry}
\usepackage{amsmath,amssymb,mathtools}
\usepackage[scr=rsfs,cal=euler]{mathalfa}
\usepackage{xfrac}
\usepackage[unicode,hidelinks,bookmarks=false]{hyperref}
\newcommand{\ie}{i.e.\ }
\newcommand{\cf}{cf.\ }
\newcommand{\resp}{respectively\ }
\newcommand{\Subsection}{\subsection}
\providecommand{\nobreakdash}{\nobreak\hbox{-}}
\setlength{\emergencystretch}{2em}
\multlinegap0pt
\usepackage[shortlabels]{enumitem}
\usepackage{amssymb}
\usepackage{mathtools}
\usepackage{tikz}
\usepackage{xcolor}
\newtheorem{proposition}{Proposition}
\newtheorem{theorem}{Theorem}
\newtheorem{corollary}{Corollary}
\newtheorem{lemma}{Lemma}
\DeclareMathOperator{\Cay}{Cay}
\newcommand{\Ghat}{\widehat{G}}
\DeclareMathOperator{\Jac}{J}
\newcommand{\Q}{\mathbb{Q}}
\newcommand{\calA}{\mathcal{A}}
\newcommand{\calO}{\mathcal{O}}
\DeclareMathOperator{\ch}{char}
\title{Spectral Algebras of Abelian Cayley Graphs}
\author{Deep Bhattacharjee, Priyabrata Mandal, Ushashi Bhattacharya}
\date{Source: arXiv:2609.09222v1 (CC BY 4.0)}
\begin{document}
\maketitle

\noindent\emph{Source and licence.} Spectral Algebras of Abelian Cayley Graphs. Version arXiv:2609.09222v1 (CC BY 4.0), distributed by arXiv under the Creative Commons Attribution 4.0 International licence, http://creativecommons.org/licenses/by/4.0/, as recorded in the arXiv licence metadata for this exact version and shown on the abstract page https://arxiv.org/abs/2609.09222v1. Full licence notice: \texttt{licenses/arxiv-2609.09222-CC-BY-4.0.txt}.\par\smallskip

\noindent\textit{Editorial scope.} Counted unit: the article's theorem thm:main (source0.tex lines 460--477). The article's complete original proof of it is delivered with it (source0.tex lines 479--507, one \texttt{proof} environment of 29 lines). No mathematics is written, repaired or re-derived: the statement and the proof are the article's own lines, in the article's own notation, and no retained line is substituted or re-flowed.  Retained uncounted as the source-local context the proof needs: lines 280--289; lines 314--340; lines 365--370; lines 444--448.  Nothing was omitted from any retained span, so the package carries no omission record.  No retained line contains a citation, so the wrapper carries no bibliography and no bibliography entry.  No retained line contains a figure environment, an image-inclusion command or drawing code, so no image is delivered and the package declares no asset dependency.  Editorial formatting only, in addition: an AMS \texttt{amsart} preamble that restates the article's own declarations and macros used by the retained lines, namely the environments \texttt{proposition}, \texttt{theorem}, \texttt{corollary}, \texttt{lemma} on separate counters, together with the standard packages those lines need.  The retained environments print in this document as Proposition 1, Theorem 1, Corollary 1, Lemma 1 in order of appearance, whereas the article prints the same items with the numbering of its own full document.  The complete native source of the article as distributed in the arXiv e-print for this exact version is bundled unchanged as \texttt{sources/209793/source0.tex}.
\begin{proposition}[Quotient-ring presentation]
\label{prop:quotient}
Let $m = m_{\Gamma,K} \in K[x]$ be the minimal polynomial of $A$ over $K$.
The evaluation map $\mathrm{ev}_A : f \mapsto f(A)$ induces a $K$-algebra
isomorphism
\[
  \calA_K(\Gamma) \;\cong\; K[x]/(m_{\Gamma,K}),
\]
and $\dim_K \calA_K(\Gamma) = \deg m_{\Gamma,K}$.
\end{proposition}

\begin{theorem}[Abelian DFT diagonalisation]
\label{thm:dft}
Let $G = \{g_0, \ldots, g_{N-1}\}$, $\Ghat = \{\chi_0, \ldots, \chi_{N-1}\}$,
and $K$ a field with $\ch(K) \nmid N$.

Since $\ch(K) \nmid N$, the polynomial $x^N - 1$ is separable over $K$
(its derivative $N x^{N-1}$ is non-zero and coprime to $x^N - 1$).
The algebraic closure $\overline{K}$ therefore contains exactly $N$ distinct
$N$-th roots of unity; fix a primitive one, $\zeta_N \in \overline{K}$.
Regard each character $\chi_j \in \Ghat$ as a group homomorphism
$G \to \overline{K}^*$ whose values are $N$-th roots of unity in $\overline{K}$,
and define the \emph{DFT matrix}
$F = (\chi_j(g_i))_{0 \le i,j < N} \in M_N(\overline{K})$.

Then $F$ is invertible with
$F^{-1} = N^{-1}(\chi_j(-g_i))^\top$, and
\begin{equation}
\label{eq:Fdiag}
  F^{-1} A F = \mathrm{diag}(\lambda_{\chi_0}, \ldots, \lambda_{\chi_{N-1}})
  \quad \text{in } M_N(\overline{K}),
\end{equation}
where $\lambda_\chi = \sum_{s \in S} \chi(s)$.
Consequently, $A$ is diagonalisable over $\overline{K}$, and its minimal
polynomial $m_{G,S,K} \in K[x]$ is squarefree over $K$: since $A$ has $N$
distinct eigenvalues in $\overline{K}$, no irreducible factor of $m_{G,S,K}$
over $K$ can occur with multiplicity greater than one.
\end{theorem}

\begin{corollary}[Cyclotomic containment]
\label{cor:cyclo}
For $K = \Q$ and $S = -S$, every eigenvalue $\lambda_\chi$ lies in the
maximal real subfield $\Q(\zeta_N)^+ = \Q(\zeta_N + \zeta_N^{-1})$.
In particular, $[\Q(\lambda_\chi) : \Q]$ divides $\varphi(N)/2$.
\end{corollary}

\begin{lemma}[Coprimality]
\label{lem:coprime}
For distinct orbits $\calO \ne \calO'$,
$\gcd(\Psi_{\calO,K}, \Psi_{\calO',K}) = 1$ in $K[x]$.
\end{lemma}

\begin{theorem}[Character-orbit decomposition]
\label{thm:main}
Let $G$ be a finite abelian group of order $N$, $S \subseteq G
\setminus \{0\}$ symmetric, and $K$ a field with $\ch(K) \nmid N$.  Then:
\begin{enumerate}[label=\textup{(\roman*)}]
  \item $m_{G,S,K}(x) = \prod_{\calO \in \calO_K(G,S)} \Psi_{\calO,K}(x)$,
        a squarefree product of distinct monic irreducibles over $K$;
  \item there is a $K$-algebra isomorphism
        \[
          \calA_K(\Cay(G,S))
          \;\cong\;
          \prod_{\calO \in \calO_K(G,S)} K[x]/(\Psi_{\calO,K}(x));
        \]
  \item $\calA_K(\Cay(G,S))$ is semisimple;
  \item for $K = \Q$, every factor $\Q[x]/(\Psi_{\calO,\Q})$ is a number
        field contained in $\Q(\zeta_N)^+$.
\end{enumerate}
\end{theorem}

\begin{proof}
\emph{Step 1: squarefreeness.}  Since $\ch(K) \nmid N$, the polynomial
$x^N - 1$ is separable over $K$, and $A$ is diagonalisable over
$\overline{K}$ (Theorem~\ref{thm:dft}).  A diagonalisable matrix has
squarefree minimal polynomial; equivalently, every irreducible factor of
$m_{G,S,K}$ over $K$ is separable (no repeated roots in $\overline{K}$),
and no irreducible can appear with multiplicity $\ge 2$ in $m_{G,S,K}$
without contradicting diagonalisability.

\emph{Step 2: orbit factorisation.}  The eigenvalues of $A$ over
$\overline{K}$ are exactly $\{\lambda_\chi : \chi \in \Ghat\}$.  The root
set of $m_{G,S,K}$ is Galois-stable, so its irreducible factors over $K$
are precisely $\{\Psi_{\calO,K} : \calO \in \calO_K(G,S)\}$, one per orbit.
This gives part~(i).

\emph{Step 3: Wedderburn.}  Proposition~\ref{prop:quotient} gives
\[
\calA_K(\Cay(G,S)) \cong K[x]/(m_{G,S,K})
= K[x]/\!\bigl(\prod_\calO \Psi_{\calO,K}\bigr).
\]
The factors are pairwise coprime
(Lemma~\ref{lem:coprime}), so the Chinese Remainder Theorem gives part~(ii).
Each factor $K[x]/(\Psi_{\calO,K})$ is a field.

\emph{Step 4: semisimplicity and cyclotomic containment.}  A product of
fields has $\Jac = 0$, giving~(iii).  For~(iv), Corollary~\ref{cor:cyclo}
gives $\lambda_\calO \in \Q(\zeta_N)^+$ for every orbit representative, so
$\Q[x]/(\Psi_{\calO,\Q}) \cong \Q(\lambda_\calO) \subseteq \Q(\zeta_N)^+$.
\end{proof}

\end{document}
